\documentclass[11pt,a4paper]{article}

%================================================
% ENCODAGE ET LANGUE
%================================================
\usepackage[utf8]{inputenc}
\usepackage[T1]{fontenc}
\usepackage[french]{babel}

%================================================
% MISE EN PAGE
%================================================
\usepackage{geometry}
\geometry{margin=1.55cm}

%================================================
% PACKAGE GTOLI
%================================================
\usepackage{../styles/gtoli}

%================================================
% INFORMATIONS DU DOCUMENT
%================================================
\title{\textbf{Biologie}\\[0.2cm]
\large Titre du chapitre}

\author{GToli}
\date{\today}

%================================================
% DOCUMENT
%================================================
\begin{document}

\maketitle

%================================================
\section{Notion biologique concernée}
%================================================

\begin{cadreobjectif}
Dans ce chapitre, l’élève doit être capable de :
\begin{itemize}
    \item identifier la notion biologique étudiée ;
    \item utiliser correctement le vocabulaire scientifique ;
    \item décrire un phénomène biologique ;
    \item expliquer les relations entre structures et fonctions ;
    \item interpréter un schéma, un document ou une expérience ;
    \item rédiger une réponse scientifique structurée.
\end{itemize}
\end{cadreobjectif}

%================================================
\section{Situation de départ}
%================================================

\begin{cadrebleu}{Problématique biologique}
Cette section permet d’introduire la situation biologique étudiée.

\medskip

On précisera :
\begin{itemize}
    \item le phénomène observé ;
    \item l’organisme ou le système concerné ;
    \item la question scientifique posée ;
    \item les hypothèses initiales éventuelles.
\end{itemize}
\end{cadrebleu}

%================================================
\section{Vocabulaire scientifique}
%================================================

\begin{cadretheorie}
\begin{itemize}
    \item \textbf{Terme 1 :} définition précise.
    \item \textbf{Terme 2 :} définition précise.
    \item \textbf{Terme 3 :} définition précise.
    \item \textbf{Terme 4 :} définition précise.
\end{itemize}
\end{cadretheorie}

%================================================
\section{Rappel théorique}
%================================================

\begin{cadretheorie}
Cette section contient :
\begin{itemize}
    \item les définitions essentielles ;
    \item les structures biologiques ;
    \item les fonctions associées ;
    \item les mécanismes ;
    \item les relations de cause à effet ;
    \item les niveaux d’organisation du vivant.
\end{itemize}

\medskip

\textbf{Point essentiel :} en biologie, il faut souvent relier une structure, une fonction et un mécanisme.
\end{cadretheorie}

%================================================
\section{Structure et fonction}
%================================================

\begin{cadregris}{Tableau structure -- fonction}
\renewcommand{\arraystretch}{1.35}
\begin{tabularx}{\textwidth}{>{\bfseries}p{3.5cm}X X}
\toprule
Structure & Fonction & Rôle dans le phénomène \\
\midrule
Structure 1 & Fonction à préciser & Rôle biologique \\
Structure 2 & Fonction à préciser & Rôle biologique \\
Structure 3 & Fonction à préciser & Rôle biologique \\
\bottomrule
\end{tabularx}
\end{cadregris}

%================================================
\section{Schéma biologique}
%================================================

\begin{cadregris}{Représentation du phénomène}
Insérer ici un schéma, une cellule, un organe, un cycle biologique, un arbre de transmission ou un processus.

\medskip

\begin{center}
\begin{tikzpicture}
\draw[thick] (0,0) circle (1.2);
\draw[fill=gray!20] (0,0) circle (0.45);
\node at (0,-1.7) {Cellule};
\node at (0,0) {N};
\draw[->,thick] (1.4,0) -- (3,0) node[right] {processus};
\end{tikzpicture}
\end{center}
\end{cadregris}

%================================================
\section{Mécanisme biologique}
%================================================

\begin{cadremethode}
Pour expliquer un mécanisme biologique :

\begin{enumerate}
    \item identifier la structure concernée ;
    \item préciser sa fonction ;
    \item décrire les étapes du phénomène ;
    \item relier les causes et les conséquences ;
    \item utiliser le vocabulaire scientifique ;
    \item rédiger une explication structurée.
\end{enumerate}
\end{cadremethode}

%================================================
\section{Exemple expliqué}
%================================================

\begin{cadreexemple}
\textbf{Situation :}

Insérer ici une situation biologique représentative du chapitre.

\medskip

\textbf{Analyse :}
\begin{itemize}
    \item Quel phénomène est observé ?
    \item Quelle structure biologique intervient ?
    \item Quelle fonction est mobilisée ?
    \item Quel mécanisme permet d’expliquer la situation ?
\end{itemize}

\medskip

\textbf{Explication :}

Présenter ici l’explication complète du phénomène.
\end{cadreexemple}

%================================================
\section{Lecture de document}
%================================================

\begin{cadreenonce}
\textbf{Document :}

Insérer ici un texte, un graphique, une photographie, un schéma ou un tableau de données.

\medskip

\textbf{Question :} analyser le document et répondre en utilisant le vocabulaire scientifique.
\end{cadreenonce}

\begin{cadreaide}
Pour analyser un document en biologie :
\begin{itemize}
    \item identifier le type de document ;
    \item repérer les informations importantes ;
    \item relier le document à la théorie ;
    \item interpréter les données ;
    \item rédiger une réponse claire.
\end{itemize}
\end{cadreaide}

%================================================
\section{Expérience ou observation}
%================================================

\begin{cadreactivite}
Cette section peut contenir :
\begin{itemize}
    \item une observation microscopique ;
    \item une expérience ;
    \item un protocole expérimental ;
    \item un tableau de résultats ;
    \item une interprétation scientifique ;
    \item une conclusion.
\end{itemize}
\end{cadreactivite}

%================================================
\section{Correction détaillée}
%================================================

\begin{cadresolution}
La correction doit contenir :
\begin{enumerate}
    \item l’identification du phénomène ;
    \item les structures concernées ;
    \item les fonctions associées ;
    \item l’explication du mécanisme ;
    \item l’interprétation du document ou de l’expérience ;
    \item la conclusion scientifique.
\end{enumerate}
\end{cadresolution}

%================================================
\section{Erreurs fréquentes}
%================================================

\begin{cadreattention}
Les erreurs fréquentes en biologie sont :
\begin{itemize}
    \item confondre structure et fonction ;
    \item utiliser un vocabulaire imprécis ;
    \item décrire sans expliquer ;
    \item oublier les relations de cause à effet ;
    \item mal interpréter un schéma ;
    \item répondre sans s’appuyer sur le document.
\end{itemize}
\end{cadreattention}

%================================================
\section{Exercices progressifs}
%================================================

\begin{cadreactivite}
\begin{enumerate}
    \item définir les mots-clés ;
    \item compléter un schéma ;
    \item associer structure et fonction ;
    \item expliquer un mécanisme ;
    \item analyser un document ;
    \item rédiger une conclusion scientifique.
\end{enumerate}
\end{cadreactivite}

%================================================
\section{Synthèse}
%================================================

\begin{cadresynthese}
À retenir :

\begin{itemize}
    \item phénomène étudié : \dotfill
    \item structure importante : \dotfill
    \item fonction principale : \dotfill
    \item mécanisme : \dotfill
    \item vocabulaire essentiel : \dotfill
    \item erreur à éviter : \dotfill
\end{itemize}
\end{cadresynthese}

%================================================
\section{Auto-évaluation}
%================================================

\begin{cadregris}{Je vérifie ma maîtrise}
\begin{itemize}
    \item[$\square$] Je connais le vocabulaire scientifique.
    \item[$\square$] Je sais reconnaître les structures.
    \item[$\square$] Je sais relier structure et fonction.
    \item[$\square$] Je sais expliquer un mécanisme.
    \item[$\square$] Je sais analyser un document.
    \item[$\square$] Je sais rédiger une réponse scientifique.
\end{itemize}
\end{cadregris}

\end{document}